%% ============================================================
%% ICI2ST_LegalShadow_Coverletter
%% Engineering Proceedings (MDPI, ISSN 2673-4591) · ICI2ST 2026 volume
%% ============================================================
\documentclass[11pt,a4paper]{article}

\usepackage[T1]{fontenc}
\usepackage[utf8]{inputenc}
\usepackage{lmodern}
\usepackage[english]{babel}
\usepackage[left=2.5cm,right=2.5cm,top=2.5cm,bottom=2.5cm]{geometry}
\usepackage{microtype}
\usepackage{parskip}
\usepackage{enumitem}
\usepackage{xcolor}
\usepackage[hidelinks]{hyperref}

\setlist[itemize]{leftmargin=1.4em, itemsep=2pt, topsep=4pt}
\newcommand{\Psiglob}{$\Psi_{\mathrm{global}}$}
\pagestyle{empty}

\begin{document}

%% ── Sender block ────────────────────────────────────────────
{\raggedright
\textbf{Patricio M. Paccha-Angamarca}\\
Software Engineering Creation, Application and Management Group (GI-CreAGSoft)\\
National Polytechnic School, Ladrón de Guevara E11-253, Quito 170525, Ecuador\\
\href{mailto:patricio.paccha@epn.edu.ec}{patricio.paccha@epn.edu.ec} \quad ORCID: 0000-0001-6285-7390\\[6pt]
Quito, Ecuador --- \today\par}

\vspace{10pt}

%% ── Addressee ───────────────────────────────────────────────
{\raggedright
\textbf{To the Guest Editors}\\
ICI2ST 2026 --- International Conference on Information Systems and Software Technologies\\
\emph{Engineering Proceedings} (MDPI, ISSN 2673-4591)\\
\emph{cc:} Editorial Office, \emph{Engineering Proceedings}\par}

\vspace{12pt}

\textbf{Subject:} Submission of the proceeding paper \emph{``LegalShadow: An Ontology-Guided Digital Shadow Architecture for Semantic Validation of Judicial Digital Ecosystems''}

\vspace{10pt}

Dear Guest Editors,

We are pleased to submit the manuscript cited above for consideration as a Proceeding Paper in the ICI2ST 2026 volume of \emph{Engineering Proceedings}. The work was prepared for presentation at ICI2ST 2026, hosted by the National Polytechnic School in Quito, and has not been published elsewhere nor is it under consideration by any other venue.

\vspace{4pt}
\textbf{Relevance to the scope of \emph{Engineering Proceedings}.}
\emph{Engineering Proceedings} provides a forum for peer-reviewed engineering research presented at academic conferences, and its scope explicitly covers \emph{computer and software engineering}. Our contribution sits squarely in that area, and in three specific senses.

\begin{itemize}
  \item \textbf{It is a software architecture paper, not a policy paper.} The core artefact is a containerised four-tier system---acquisition, semantic, validation and presentation---whose design decisions (unidirectional data flow, container-level isolation between evidence harvesting and scoring, a parameterised acquisition target) are argued from engineering requirements and reported so that they can be reproduced and criticised as engineering.
  \item \textbf{It transfers a cyber-physical systems abstraction to a new class of system.} We adopt Kritzinger et al.'s digital model / shadow / twin distinction and instantiate the \emph{digital shadow} configuration---automatic real-to-virtual flow, human-controlled feedback---over a public institutional software ecosystem. This is the engineering rationale for the design: acting automatically on public infrastructure would be improper, so the artefact is deliberately built to observe and diagnose rather than to actuate.
  \item \textbf{It addresses an open verification problem in the digital twin literature.} A recent systematic review reports no established method to verify that a twin actually materialises its semantic model. We contribute a hierarchical coverage function \Psiglob{} and a gap function $\delta$ that make this verifiable: both are bounded, monotone, computable in linear time in the scoring stage, and interpretable dimension by dimension.
\end{itemize}

The application domain is the judicial sector, but the object of study is a software ecosystem and the contribution is an instrument for assessing it. We believe this is a good fit for a journal whose remit is engineering work presented at conferences, and for an ICI2ST volume centred on information systems and software technologies.

\vspace{4pt}
\textbf{Contribution and main result.}
The paper makes three contributions: \textbf{(C1)} the LegalShadow reference architecture and its coverage/gap model $\langle\Psi,\delta\rangle$; \textbf{(C2)} a non-invasive acquisition protocol with per-snapshot cryptographic provenance, specified as an executable procedure; and \textbf{(C3)} an empirical study of the Ecuadorian Judiciary Council over 17 dated runs. The headline finding is a measurable distance between declared and machine-actionable governance: the institution scores 46/100 on the four dimensions it instruments itself, yet its observable digital surface materialises only \Psiglob{}~$=34.9\%$ of the semantics that the ten frameworks it invokes prescribe, with three dimensions at zero anchored concepts. The finding is also reflexive, and this is what we consider most interesting for a software engineering readership: the observed institution lacks precisely the embedded structured data that the shadow itself needs in order to represent it.

\vspace{4pt}
\textbf{Reproducibility.}
Every figure reported in the paper is recomputable by a third party. The acquisition protocol is read-only and identified, honours the target's \texttt{robots.txt}, and records each capture with a UTC timestamp and SHA-256 checksum in an append-only hash-chained ledger; re-hashing the 170 stored snapshots yields 170 matches and 0 mismatches. The source code, the OWL governance ontology, the dated JSON-LD shadows and the evidence manifests are openly deposited, and the paper reports the ledger head so that any reviewer can verify that no cited snapshot was altered. We would welcome reviewers exercising the artefact directly: the environment builds with a single \texttt{docker compose} invocation.

\vspace{4pt}
\textbf{Ethical and legal considerations.}
All observed data originate from unauthenticated public resources. The agent identifies itself with an explicit \emph{User-Agent} stating its purpose, submits no forms, performs no writes, keeps at most one request in flight per target and caps each response, so that the measured load on the observed portal is negligible (median latency 788.5~ms over 187 requests). No personal data were collected, and the study makes no claim about individuals.

\vspace{4pt}
\textbf{Declarations.}
The manuscript is original, has been read and approved by all authors, and reports no external funding. The authors declare no conflicts of interest and no financial or personal relationship with the observed institution. All prior work by the authors on which this study builds is cited explicitly in the introduction. We confirm that the paper conforms to the \emph{Engineering Proceedings} template and author instructions, and we are willing to suggest independent reviewers should the Guest Editors find it useful.

\vspace{4pt}
Thank you for considering our submission. We look forward to the reviewers' comments and remain at your disposal for any clarification.

\vspace{14pt}

{\raggedright
Yours sincerely,\\[22pt]
\textbf{Patricio M. Paccha-Angamarca}\\
\emph{on behalf of the co-authors} Erwin J. Sacoto-Cabrera and Víctor V. Velepucha-Bonett\\
Corresponding author --- \href{mailto:patricio.paccha@epn.edu.ec}{patricio.paccha@epn.edu.ec}\par}

\end{document}
